%
% File hicss.tex
%
% Contact: Holm Smidt, hicss@hawaii.edu
%%
%%
%% Based on the style files for ACL 2015 by
%% car@ir.hit.edu.cn, gdzhou@suda.edu.cn


\documentclass[10pt]{article}
\input{hicss-packages.tex}
\newcommand{\sansserifformat}[1]{\fontfamily{cmss}{ #1}}%

% ---- Packages needed for the Fig.1/Fig.2 TikZ diagrams (see paper/ .tex) ----
% If any of these are already pulled in by hicss-packages.tex, the re-load is
% harmless. amssymb is used for math symbols in the method section.
\usepackage{amssymb}
\usepackage{tikz}
\usetikzlibrary{arrows.meta,positioning,shapes.geometric,fit,backgrounds,calc}

% Convenience macro for our method's name.
\newcommand{\method}{UA-ProtoPNet}

\setlength\titlebox{7cm}

% You can expand the titlebox if you need extra space
% to show all the authors. Please do not make the titlebox
% smaller than 5cm (the original size).
% 5cm is the default for a single row list of authors


% title for manuscript without author info
\title{Explaining Why, Not Just How Much: Competing-Prototype Uncertainty
for Human--AI Collaboration in Skin-Lesion Screening}

% title for final manuscript
% comment for initial manuscript
% \title{...}

% Comment this for initial manuscript
% Uncomment this for final manuscript
 \author{Himson Chapagain \\
  Texas State University \\
  {\underline{rnx16@txstate.edu}} \\ }

\date{}

\begin{document}
\maketitle

\begin{abstract}
When a clinical AI system is uncertain, the quality of the hand-off to a human
expert determines patient safety outcomes. Yet most systems hand the clinician
a label and, at best, a scalar uncertainty score: they convey \emph{that} a
model is unsure but not \emph{why}, offering no basis for action at the point
of hand-off. We present \method{}, an interpretable classifier that explains
its uncertainty by surfacing the competing prototypical exemplars driving an
ambiguous decision---when torn between melanoma and a benign nevus, it shows a
real, provenance-labeled example of each. The explanation is grounded in the
model's actual decision: the top competing classes by logit and, for each, the
prototype contributing most to that class. On HAM10000 (dermoscopy), \method{}
matches standard ProtoPNet accuracy while uniquely providing case-level
explanations of its uncertainty. On LIDC/LUNA22 we further show, across random
seeds, that no method---including Monte-Carlo Dropout and Deep
Ensembles---predicts inter-radiologist disagreement from images, arguing that
diagnostic ambiguity should be \emph{surfaced} to the clinician rather than
silently resolved.
\end{abstract}

\subsubsection*{Keywords:}

Explainable AI; uncertainty quantification; human--AI collaboration; clinical
decision support; medical image classification.


\section{Introduction}

Artificial intelligence is moving from research benchmarks into the front line
clinical triage: dermoscopy classifiers screen skin lesions~\cite{esteva2017dermatologist} and CT (Computed Tomography) models flag
lung nodules~\cite{setio2017luna}, often as a first pass before a specialist. In these settings, the
unit of value is not a standalone prediction but a \emph{collaboration}: the
model proposes, and a clinician might accept or deny. The moment that matters most is
therefore the \emph{hand-off}---when the model is unsure and must consult/enlist the
human. When that hand-off procedure is a black-box deferral, patient safety is at risk: a clinician
handed an opaque or black-box deferral has no actionable interpretation or explanation to make the final decision.
In our work, we redesign that hand-off procedure.

Most deployed systems handle this moment poorly. They output a class label and,
if anything, a scalar uncertainty or confidence number. A score of ``$U = 0.84$''
tells a clinician \emph{that} the model hesitated, but not \emph{about what}, or
\emph{why}. It is an alarm without an explanation---a deferral that enlists the
human without helping them. Figure~\ref{fig:workflow} presents the
collaboration workflow this paper proposes: one in which uncertain cases are
escalated not with a score, but with the concrete competing evidence driving the
model's doubt.

Previously, two research works,  each addressed half of the problem. Uncertainty
quantification methods---Monte-Carlo (MC) Dropout
\parencite{gal2016dropout} and Deep Ensembles
\parencite{lakshminarayanan2017ensembles}---estimate \emph{how much} a model is
unsure, but as opaque scalars: they cannot say which classes or features drive
the doubt. Interpretable models such as ProtoPNet
\parencite{chen2019protopnet} explain a \emph{prediction} by pointing to
prototypical training exemplars (``this looks like that''), but they explain the
chosen class, not the \emph{uncertainty} between classes. Neither thread
delivers what a clinical hand-off needs: an explanation of the ambiguity itself.

We close this gap with \method{} (Uncertainty-Aware ProtoPNet). When the model
is torn, it produces a \emph{competing-prototype explanation}: it identifies the
two classes the decision is actually split between and, for each, retrieves the
real training exemplar most responsible for that class's score. For an uncertain
pigmented lesion, this means showing a histopathology-confirmed melanoma
exemplar \emph{and} a benign nevus exemplar side by side (Figure~\ref{fig:workflow}), alongside the model's
class-probability distribution and an uncertainty score
. The clinician sees not just that the model is
unsure, but the concrete visual evidence pulling it in each direction---turning
an opaque/black-box deferral into a well-explained hand-off.

Crucially, the explanation is \emph{grounded in the decision}. A naive
prototype-similarity ranking surfaces clinically irrelevant ``magnet''
prototypes, because raw per-class similarities are near-uniform; the class
separation lives in the linear decision layer. We therefore derive the competing
prototypes from the model's logits and their prototype contributions, therefore, the
explanation always reflects the classes the model is genuinely weighing.

Our contribution is three-folds:
\begin{itemize}
\setlength\itemsep{0.1em}
\item A \textbf{clinical hand-off workflow design} for uncertainty-aware
human--AI collaboration in medical image triage: an interpretable,
uncertainty-aware system that routes confident cases to automated triage and
escalates uncertain cases to the clinician with actionable evidence, not just a
score (Figures~\ref{fig:arch} and~\ref{fig:workflow}).
\item A \textbf{decision-grounded competing-prototype explanation mechanism}
that makes the hand-off actionable: when the model is torn, it surfaces the two
real, provenance-labeled training exemplars driving the ambiguity---turning an
opaque deferral into an informed clinical judgment
(Section~\ref{sec:results}, Figure~\ref{fig:hero}).
\item A rigorous \textbf{empirical finding that aleatoric uncertainty is not
image-predictable}: no evaluated method---including MC-Dropout and Deep
Ensembles---can predict inter-radiologist disagreement from lung nodule images,
arguing that \emph{surfacing} ambiguity to clinicians is the only responsible
design choice rather than attempting to resolve it
(Figure~\ref{fig:scatter}, Table~\ref{tab:negative}).
\end{itemize}


\section{Related Work}

\textbf{AI for medical-image triage.} Deep classifiers approach specialist-level
performance on dermoscopy \parencite{esteva2017dermatologist,
tschandl2018ham10000} and lung-nodule analysis \parencite{armato2011lidc,
setio2017luna}. As these tools enter screening workflows, attention has shifted
from raw accuracy to how they integrate with human decision-making
\parencite{topol2019high}.

\textbf{Uncertainty quantification.} MC-Dropout \parencite{gal2016dropout} and
Deep Ensembles \parencite{lakshminarayanan2017ensembles} are the standard tools
for estimating predictive uncertainty, and calibration is commonly assessed via
the Expected Calibration Error \parencite{guo2017calibration}. In medicine,
uncertainty is argued to be essential for safe deployment
\parencite{begoli2019need, kompa2021second}. These methods, however, yield a
scalar; they do not explain the source of the uncertainty.

\textbf{Interpretable and prototype models.} Rather than post-hoc saliency,
\textcite{rudin2019stop} argues for models interpretable by construction.
ProtoPNet \parencite{chen2019protopnet} classifies by comparing image regions to
learned prototypes that are projected onto real training patches. Prototype models explain the predicted class; we extend them to explain the
\emph{uncertainty} between classes. Recent work has argued for jointly
considering explainability and uncertainty in medical AI
\parencite{Salvi2025uncertainty}; our approach instantiates this idea
with prototype-based explanations grounded in competing decision evidence.

\textbf{Human--AI collaboration and deferral.} A growing literature studies when and how AI should defer to humans
\parencite{mozannar2020consistent} and how explanations shape trust and reliance
\parencite{ghassemi2021false}. A central concern in this literature is
\emph{appropriate reliance}---calibrating human trust so clinicians accept
AI recommendations when the model is reliable and scrutinize them when it is
not \parencite{ghassemi2021false,mozannar2020consistent}. Existing deferral systems, however, treat the hand-off as a binary routing
decision; they do not specify what information the human receives at the moment
of escalation, even in complex settings such as AI-supported sepsis diagnosis in
the ICU \parencite{Zhang2024rethinking}. We contribute to this literature by
operationalizing the hand-off itself as a first-class design object:
\method{} determines not only \emph{when} to defer, but \emph{what} to surface
to the clinician at the point of deferral.


\section{Methodology}
\label{sec:method}

% =========================== Figure 2 (architecture) =========================
% TO USE THE REAL FIGURE: replace the \includegraphics line below with
%   \input{fig2_arch}

\begin{figure*}[!htbp]\centering
    \resizebox{\textwidth}{!}{%
      \input{fig2_arch}
    }% <-- uncomment this and delete the line above
  \caption{\textbf{\method{} architecture.} Blue: standard ProtoPNet
  (prototypes projected onto real training patches via the push step). Orange:
  our extension deriving, from the same decision, (A) an uncertainty score and
  (B) a decision-grounded competing-prototype explanation.}
  \label{fig:arch}
\end{figure*}

\subsection{Background: ProtoPNet}

A ProtoPNet \parencite{chen2019protopnet} comprises a convolutional backbone
$f$, a set of $C{\times}K$ prototype vectors ($K$ per class for $C$ classes), a
prototype similarity layer, and a linear classification layer. The backbone maps
an input $x$ to a feature map; each prototype $p_{c,k}$ is compared to all
spatial patches, and a global max-pool yields a single activation
$a_{c,k}=\max_{\text{patches}} g(p_{c,k}, \cdot)$ measuring how strongly the
strongest patch resembles that prototype, where $g$ converts an $L_2$ distance to
a similarity. The activation vector $a\in\mathbb{R}^{CK}$ is mapped by a linear
layer $W$ to class logits $z = Wa$ and probabilities $p=\mathrm{softmax}(z)$.
A \emph{prototype push} step projects each prototype onto its nearest training
patch, so every prototype is a viewable, real exemplar.

\subsection{Decision-grounded Competing Prototypes}

For any given input, \method{} interprets the model's \emph{uncertainty} by identifying
conflicting classes and the exemplars driving each. A natural but
flawed approach ranks classes by their max prototype similarity (``class
evidence''); we find these are near-uniform across classes, so the top
competitors are arbitrary high-similarity prototypes unrelated to the decision.
Instead, we ground the explanation in the decision itself. Let
$\hat{c}_1, \hat{c}_2$ be the top-two classes by logit $z$. For each competing
class $c\in\{\hat{c}_1,\hat{c}_2\}$ we select the prototype contributing most to
that class's logit,
\begin{equation}
k^\ast(c) \;=\; \operatorname*{arg\,max}_{k} \; a_{c,k}\, W_{c,k},
\end{equation}
and retrieve its pushed training patch (with its metadata, e.g.\ diagnosis,
confirmation type, and body site). The result is a pair of real, provenance-labeled exemplars---one per competing class---which explains \emph{why} the model
is uncertain (Figure~\ref{fig:arch}, branch~B).

\subsection{Uncertainty Score}

We quantify uncertainty as the normalized Shannon entropy of the predictive
distribution,
\begin{equation}
U \;=\; -\frac{1}{\log C}\sum_{c=1}^{C} p_c \log p_c \;\in\; [0,1],
\end{equation}
with an optional temperature $T$ applied to the logits and calibrated on a
validation split. Using the predictive distribution (rather than raw prototype
similarities) ensures the scalar uncertainty and the competing-prototype
explanation describe the \emph{same} decision. The model thus outputs, for every
case, a prediction $\hat{y}$, an uncertainty $U$, and---when $U$ is high---the
two competing exemplars (Figure~\ref{fig:workflow}).


\section{Experimental Setup}
\label{sec:setup}

% =========================== Figure 1 (workflow) =============================
% TO USE THE REAL FIGURE: replace the \includegraphics line below with
%   \input{fig1_workflow}
\begin{figure*}[!htbp]\centering
    \resizebox{\textwidth}{!}{%
      \input{fig1_workflow}
    }% <-- uncomment this and delete the line above
  \caption{\textbf{Human--AI collaboration workflow for uncertainty-aware
  skin-lesion screening.} A conventional scalar-uncertainty model (top) reports
  only \emph{that} it is uncertain; \method{} (bottom) routes confident cases to
  automated triage, whereas on uncertain cases, produces a competing-prototype
  explanation that turns an opaque deferral into an actionable clinician
  hand-off.}
  \label{fig:workflow}
\end{figure*}

\textbf{Datasets.} \emph{HAM10000} \parencite{tschandl2018ham10000} contains
10{,}015 dermoscopic images across seven diagnostic classes (melanoma,
melanocytic nevus, basal cell carcinoma, actinic keratosis, benign keratosis,
dermatofibroma, vascular lesion). We split $70/15/15$ stratified by class and,
critically, \emph{by lesion}: because multiple images can depict the same
physical lesion, we assign each lesion entirely to one split to prevent
train/test leakage that would inflate accuracy and calibration. Images are
resized to $224{\times}224$ with ImageNet normalization and standard
augmentation.

\emph{LIDC/LUNA22} \parencite{armato2011lidc, venkadesh2022luna22} provides lung-nodule
CT patches, each rated for malignancy on a 1--5 scale by up to four radiologists.
The \emph{standard deviation} of these ratings is a ground-truth signal of
inter-radiologist disagreement. We extract the centered axial slice of each
nodule and crop to the nodule region (we found this crop essential: without it
the nodule is a small fraction of the field of view and the model fails to
learn). We study a binary benign-vs-malignant task on the confident nodules
(mean rating $\le 2.5$ or $\ge 3.5$), holding out the indeterminate nodules
(mean rating in $(2.5,3.5)$) as an ambiguous evaluation set. Single-channel
patches are replicated to three channels for the ImageNet-pretrained backbone.

\textbf{Baselines.} We compare \method{} against (i) a vanilla ProtoPNet, (ii)
MC-Dropout \parencite{gal2016dropout} ($30$ stochastic passes; variance as
uncertainty), and (iii) a Deep Ensemble \parencite{lakshminarayanan2017ensembles}
(independently trained members; variance as uncertainty). All share the backbone
and prototype configuration; only \method{} provides explanations.

\textbf{Metrics.} We report classification accuracy; Expected Calibration Error
(ECE, $\downarrow$) \parencite{guo2017calibration}; the AUROC of uncertainty for
error detection ($\uparrow$); and, for LIDC/LUNA22, the Pearson correlation
between model uncertainty and the radiologist-rating standard deviation, plus an
image-independent \emph{difficulty} correlation (uncertainty vs.\ proximity of
the mean malignancy to the decision boundary), which the model never observes
during training.

\textbf{Implementation.} The backbone is an ImageNet-pretrained ResNet-50 with
$K{=}10$ prototypes per class. Training follows the ProtoPNet schedule
(warm-up, joint optimization, prototype push, and final-layer fine-tuning). All
metrics are reported as mean$\pm$std over multiple random seeds, with the data
split held fixed so that only model initialization and training vary.


\section{Results}
\label{sec:results}

% =========================== Table 1 (main comparison) =======================
\begin{table*}[!htbp]\centering
\caption{Classification and uncertainty metrics on HAM10000 (mean$\pm$std over
three seeds). \method{} matches the accuracy of standard ProtoPNet while being
the only method that provides a prototype-grounded explanation of its
uncertainty (final row). Variance-based baselines (MC-Dropout, Ensemble) are
better calibrated; \method{} and vanilla share a deterministic softmax. Best
value per metric in \textbf{bold}.}
\label{tab:main}
\begin{tabular}{l cccc}
\hline
\textbf{Metric} & \textbf{Vanilla} & \textbf{MC Dropout} & \textbf{Ensemble}
  & \textbf{\method{} (Ours)} \\
\hline
Accuracy ($\uparrow$)          & $0.69{\pm}0.02$ & $0.69{\pm}0.02$ & $\mathbf{0.72{\pm}0.01}$ & $0.69{\pm}0.02$ \\
ECE ($\downarrow$)             & $0.25{\pm}0.01$ & $\mathbf{0.10{\pm}0.01}$ & $0.12{\pm}0.01$ & $0.25{\pm}0.02$ \\
Uncertainty AUROC ($\uparrow$) & $0.69{\pm}0.11$ & $\mathbf{0.84{\pm}0.02}$ & $0.79{\pm}0.01$ & $0.70{\pm}0.09$ \\
Prototype explanation          & No & No & No & \textbf{Yes} \\
\hline
\end{tabular}
\end{table*}

\subsection{\method{} Preserves Accuracy}

Table~\ref{tab:main} reports classification and uncertainty metrics on HAM10000
(mean$\pm$std over three seeds). Because \method{} adds an explanation layer
without altering the classifier, its accuracy matches vanilla ProtoPNet (both
$\approx 0.69$) and trails the best method---the ensemble ($0.72$)---by about
three points. The scalar uncertainty metrics, however, are not where \method{}
competes: the variance-based methods (MC-Dropout and the ensemble) are better
calibrated (ECE $0.10$--$0.12$ vs.\ $0.25$) and detect errors better
(AUROC $0.79$--$0.84$) than the single deterministic softmax shared by \method{}
and vanilla ProtoPNet. The decisive distinction is the final row: \method{} is
the only method that can \emph{explain} its uncertainty, at no accuracy cost
relative to its ProtoPNet base.

Crucially, \emph{no} scalar signal here captures genuine ambiguity. Uncertainty
does not correlate with radiologist disagreement for any method
(Section~\ref{sec:results}, Table~\ref{tab:negative}).The ``ambiguous'' classes
are dominated by nevus---common and easy for the model---so that group shows
lower uncertainty than the rare ``clear'' classes, making the contrast reflect
class frequency rather than genuine ambiguity. Rather than a weakness of our method, this is its empirical
motivation: when scalar uncertainty is uninformative or confounded, the
actionable signal for a clinician is an \emph{explanation} of the ambiguity, not
another number.

\subsection{Detailed Explanation of Uncertain Prediction Case}

Figure~\ref{fig:hero} demonstrates the scenario where the model misclassify a melanoma case as a benign case. 
The lesion is in fact a melanoma, but the model assigns the highest probability
to benign keratosis ($0.58$); rather than silently committing to that benign
label, \method{} surfaces the competing exemplars---a histopathology-confirmed
melanoma exemplar alongside the benign one---and shows that melanoma is a strong
competing explanation ($0.40$). The explanation makes this near-miss more visible and explanatory,
prompting clinician review precisely where a scalar ``uncertain'' ($U=0.55$)
would have conveyed only that the model is uncertain. This is the interpretable
hand-off of Figure~\ref{fig:workflow} instantiated on a real case.

The scope of this work is explicitly stated here. The two competing exemplars and the
class-probability distribution---which rest directly on the model's decision
(logits and prototype contributions)---are the substance of the explanation. The
matched-region boxes in the left panel are a localizing cue: as is well
documented for prototype models, the region a prototype responds to most
strongly need not be a pixel-level visual match to the prototype patch, since the
similarity is computed in the network's learned feature space rather than in raw
appearance \parencite{hoffmann2021looks}. We, therefore, read the figure as ``the
model is split between a melanoma exemplar and a benign exemplar,'' not as ``this
region looks identical to that patch,'' and present the boxes as supporting
context rather than the claim itself.

The behavior generalizes beyond the melanoma case. Figure~\ref{fig:hero2} demonstrates
another ambiguous pair---actinic keratosis, a pre-malignant lesion, versus
benign keratosis---where the model is again nearly split (probabilities $0.48$
vs.\ $0.34$). Here the two retrieved exemplars are themselves visually similar
keratotic patterns, which is exactly why the decision is ambiguous: rather than
committing to a single label, \method{} presents both readings to the clinician and allows the
clinician to weigh the pre-malignant against the benign interpretation. Across both
cases, the explanation is the pair of real, provenance-labeled exemplars the
decision is split between---something no scalar-uncertainty baseline can provide.


% =========================== Figure 3 (hero explanation) =====================
\begin{figure*}[!htbp]\centering
  \includegraphics[width=0.95\linewidth]{fig3_hero}
  \caption{\textbf{Competing-prototype explanation of an uncertain prediction
  (\method{}).} Left: input lesion with the two competing prototypes' activation
  regions (red: predicted class; blue: runner-up). Center: the two competing
  training exemplars, with class, provenance, and similarity. Right: the model's
  predicted class distribution and uncertainty $U$.}
  \label{fig:hero}
\end{figure*}
\begin{figure*}[!htbp]\centering
  \includegraphics[width=0.95\linewidth]{fig_example2}
  \caption{\textbf{A second competing-prototype explanation
  (actinic vs.\ benign keratosis).} A keratotic lesion (ground truth: benign
  keratosis) on which the model is split between actinic keratosis ($0.48$) and
  benign keratosis ($0.34$), $U=0.81$. Both competing exemplars are
  histopathology-confirmed and are visually similar keratotic patterns---which
  is precisely why the decision is ambiguous. As in Figure~\ref{fig:hero}, the
  matched-region boxes localize where each prototype responds; the two competing
  exemplars and the class distribution carry the explanation.}
  \label{fig:hero2}
\end{figure*}

\subsection{Uncertainty across Classes}

Figure~\ref{fig:dist} shows the distribution of \method{}'s uncertainty by
class. The pattern is driven by class frequency rather than visual ambiguity:
the visually ambiguous classes (melanoma, nevus, benign keratosis) are
dominated by the majority nevus class, on which the model is typically
confident, whereas the rarer, visually distinct classes attract higher
uncertainty. A class-level ``ambiguous-vs-clear'' contrast therefore does not
isolate genuine ambiguity, which is why we do not report it as an
uncertainty-quality metric and instead emphasize case-level explanations
(Figure~\ref{fig:hero}).

% =========================== Figure 4 (distribution) =========================
\begin{figure}[tb]\centering
  \includegraphics[width=0.95\linewidth]{uncertainty_distribution}
  \caption{\textbf{Distribution of uncertainty by class (HAM10000 test set).}
  Each violin shows the per-class spread of \method{}'s uncertainty $U$; the
  dashed line marks a reference median ($\approx 0.28$). Per-class uncertainty
  tracks class frequency more than visual ambiguity (Section~\ref{sec:results}).}
  \label{fig:dist}
\end{figure}

\subsection{Inter-Radiologist Disagreement is not Image-Predictable}

An ususal hypothesis is that an effective uncertainty estimator would track
\emph{aleatoric} uncertainty---the genuine ambiguity captured by radiologist
disagreement. On LIDC/LUNA22 we test this directly. Figure~\ref{fig:scatter}
plots model uncertainty against the standard deviation of radiologist malignancy
ratings, and Table~\ref{tab:negative} reports the correlation for every method.
The result is unambiguous and robust across seeds: \emph{no} method
—including MC-Dropout and Deep Ensembles, which are designed for aleatoric uncertainty
—achieves a correlation distinguishable from zero (all $|r|\!<\!0.06$), and the
same holds for the image-independent difficulty correlation. Importantly, this
is not because the classifier failed to learn the task: on the binary
benign-vs-malignant nodule problem \method{} reaches $0.71\!\pm\!0.03$ accuracy,
comparable to the ensemble ($0.72\!\pm\!0.07$), so the models clearly capture
malignancy-relevant content---they simply cannot predict inter-human annotation (inter-radiologist)
\emph{disagreement}. We interpret this not as a failure of a particular model
but as evidence that disagreement is dominated by the subjectivity of human annotation rather than
recoverable image content. This finding directly motivates our design stance:
when ambiguity cannot be \emph{predicted}, the responsible system should
\emph{report} it, with an explanation, to the human.

% =========================== Figure 5 (scatter) ==============================
\begin{figure}[tb]\centering
  \includegraphics[width=0.95\linewidth]{correlation_scatter}
  \caption{\textbf{Model uncertainty vs.\ inter-radiologist disagreement
(LIDC/LUNA22),} colored by ground-truth label. Uncertainty does not track
disagreement: the indeterminate (orange) nodules span the full uncertainty range
and the trend is flat (negligible $|r|<0.1$; Table~\ref{tab:negative}).}
  \label{fig:scatter}
\end{figure}

% =========================== Table 2 (negative result) =======================
\begin{table}[tb]\centering
\caption{Inter-radiologist disagreement is not image-predictable: Pearson
correlation (mean$\pm$std over three seeds) between each method's uncertainty and
the std of radiologists' malignancy ratings on LIDC/LUNA22. All values are
indistinguishable from zero.}
\label{tab:negative}
\begin{tabular}{l c}
\hline
\textbf{Method} & \textbf{Pearson $r$ (vs.\ radiologist std)} \\
\hline
Vanilla ProtoPNet & $-0.04{\pm}0.06$ \\
MC Dropout        & $-0.05{\pm}0.04$ \\
Deep Ensemble     & $\phantom{-}0.04{\pm}0.02$ \\
\method{} (Ours)  & $-0.05{\pm}0.04$ \\
\hline
\end{tabular}
\end{table}


\section{Discussion: Implications for Human--AI Decision-Making}

\textbf{Report, do not resolve, ambiguity.} Our negative result reframes the
goal of clinical uncertainty estimation. If inter-expert disagreement cannot be
predicted from the image, then a system that reports a confident label on a
genuinely ambiguous case is misleading, and a system that reports only a scalar
``uncertain'' is unhelpful. The constructive alternative is to make ambiguity
\emph{legible}: demonstrate the human the competing evidence so they can adjudicate.
\method{} operationalizes exactly this.

\textbf{Explanations make uncertainty actionable.} A scalar uncertainty score is
difficult to act on because it is uncalibrated to the clinician's own mental
model. A pair of competing exemplars---``the model is split between this
confirmed melanoma and this benign nevus''---is immediately interpretable in
clinical terms and supports a concrete next step (biopsy, refer, or dismiss).
This aligns uncertainty-aware AI with the way clinicians already reason by
analogy and comparison.

\textbf{Selective deferral and trust calibration.} The workflow of
Figure~\ref{fig:workflow} is a selective-prediction system: confident cases are
automated, uncertain cases are escalated with an explanation. Because the
escalation carries \emph{evidence}, it supports appropriate reliance---trusting
the model when its exemplars are coherent and scrutinizing it when they are
not---rather than the blanket over- or under-reliance with opaque/black-box scoring reporting .

\textbf{Design implications.} For teledermatology and screening triage, our
results advocate (i) prioritizing interpretable, exemplar-based explanations of
uncertainty over additional scalar uncertainty machinery, and (ii) designing the
hand-off, not just the prediction, as a first-class design component of the system.


\section{Limitations and Future Work}

Our current work has several limitations. First, the explanations and
the dermatology evaluation use a single 2D modality; richer 3D context may
change the lung-nodule findings. Second, our central result on aleatoric
uncertainty is a \emph{negative} one; while robust across methods and seeds, it
establishes the absence of an image-predictable signal rather than a new
predictive capability. Third, \method{}'s quantitative advantage over baselines
is in interpretability rather than scalar metrics, on which it is competitive
but not dominant. Fourth, \method{} inherits a known property of prototype
models: the spatial region a prototype matches need not be a pixel-level visual
match to the prototype patch, because similarity is measured in the network's
learned feature space rather than in raw appearance
\parencite{hoffmann2021looks}. This is why we treat the competing exemplars and
the decision distribution as the explanation and the matched-region boxes only as
a localizing cue (Section~\ref{sec:results}); it does not affect the validity of
the competing-class account, which is grounded in the logits. Finally, we evaluate the explanations computationally and qualitatively.
A controlled clinician study is the natural next step, and its design is
straightforward: participants randomized to scalar-uncertainty versus
competing-prototype conditions, with three primary outcomes---
(1) \emph{decision accuracy} on ambiguous cases where ground truth is known,
(2) \emph{trust calibration}, measured as the correlation between stated
confidence and actual correctness across confident and uncertain model outputs,
and (3) \emph{time-to-decision}, capturing whether a richer explanation
accelerates or burdens the hand-off. The negative result on inter-radiologist
disagreement (Section~\ref{sec:results}) further motivates this study: if
ambiguity cannot be predicted from images, the value of surfacing/reporting it must
ultimately be measured in human decision outcomes.


\section{Conclusion}

We presented \method{}, an interpretable classifier that explains \emph{why} it
is uncertain by surfacing/reporting the competing prototypical exemplars driving an
ambiguous decision. The explanation is grounded in the model's actual decision,
preserves classification accuracy, and uniquely supports an informed clinical
hand-off that scalar-uncertainty baselines cannot. We further showed, rigorously,
that inter-radiologist disagreement is not recoverable from nodule images by any
evaluated method---an argument for surfacing/reporting ambiguity to clinicians rather than
silently resolving it. We advocate for interpretable, uncertainty-aware hand-off as an effective
 design pattern for human–AI collaboration in medical-image triage.


% if added before the last page, this command can help balancing columns
%\addtolength{\textheight}{-.2cm}

% ----------------------------------------------------------------------------
% References. The HICSS template uses biblatex (\printbibliography). Add the
% following keys to sample.bib (BibTeX/biblatex). Suggested sources:
%   chen2019protopnet      -- Chen et al., "This Looks Like That", NeurIPS 2019
%   gal2016dropout         -- Gal & Ghahramani, MC Dropout, ICML 2016
%   lakshminarayanan2017ensembles -- Lakshminarayanan et al., Deep Ensembles, NeurIPS 2017
%   guo2017calibration     -- Guo et al., On Calibration of Modern NNs, ICML 2017
%   tschandl2018ham10000   -- Tschandl et al., HAM10000, Scientific Data 2018
%   armato2011lidc         -- Armato et al., LIDC-IDRI, Medical Physics 2011
%   setio2017luna          -- Setio et al., LUNA16, Medical Image Analysis 2017
%   esteva2017dermatologist-- Esteva et al., Nature 2017
%   rudin2019stop          -- Rudin, Stop Explaining Black Boxes, Nature MI 2019
%   guo... / begoli2019need-- Begoli et al., Nature MI 2019
%   kompa2021second        -- Kompa et al., npj Digital Medicine 2021
%   mozannar2020consistent -- Mozannar & Sontag, Learning to Defer, ICML 2020
%   ghassemi2021false      -- Ghassemi et al., The false hope of XAI, Lancet Digital Health 2021
%   topol2019high          -- Topol, High-performance medicine, Nature Medicine 2019
% ----------------------------------------------------------------------------

\printbibliography

\end{document}
